\chapter*{Abstract}
\addcontentsline{toc}{chapter}{Abstract}

Scientific question answering in specialised domains requires both factual precision and
contextual awareness. Earth Observation (EO) is one such domain: queries combine
multiple sensors, biophysical variables, and strict spatial and temporal boundaries,
making reliable retrieval and generation particularly challenging. Open-source
Retrieval-Augmented Generation (RAG) systems struggle here: even when relevant
documents are retrieved, the assembled context fails to reflect the specific dimensions
of the question, degrading the final answer.

Existing knowledge-graph RAG approaches~\citep{elbaff2025earthobservationrag} address
retrieval but leave context curation unstructured, leading to four recurring failure
modes: delayed direct answers, missing quantitative precision, poor response structure,
and context misalignment. This thesis proposes an ontology-based three-agent pipeline
that addresses these shortcomings by explicitly modelling the semantic structure of each
query. An ontology agent parses each query into a lightweight, query-specific ontology
of attribute-value pairs across five semantic dimensions (e.g., temporal scope:
2015--2020, sensor: Sentinel-2). A refinement agent uses this ontology to filter
retrieved documents and synthesise a curated context. A generation agent then drafts and
refines the final answer.

Our experiments employ a small open-source LLM, Mistral Small 24B, across four pipeline
variants evaluated against 70 EO questions~\citep{elbaff2025earthobservationrag} using
an LLM-as-a-judge framework with Claude Sonnet~4.6~\citep{anthropic2024claude} as
judge. The ontology-enhanced pipeline achieved an overall score of 7.48, ahead of the
baseline (5.51) and a no-ontology variant (6.88). All pairwise comparisons were
significant ($p < 0.001$, $r \geq 0.77$). The architecture generalises to any
knowledge-intensive domain where queries span multiple semantic dimensions.
